\section{Introduction}

Multi-agent systems that support autonomous spawning—wherein agents may create child agents to decompose and delegate tasks—offer compelling advantages in scalability, fault isolation, and parallel computation. However, this expressiveness introduces a fundamental hazard we term \textbf{autonomous drift}: the gradual or sudden loss of oversight, accountability, and behavioral guarantees over a growing descendant population. Drift is not merely a monitoring or observability problem. It is a structural vulnerability that arises whenever the spawn graph becomes opaque to the originating system. Once an agent can spawn without the originating context retaining a verifiable reference to the descendant, the descendant's behavior escapes the safety boundary of the parent.

In existing frameworks, spawn graphs become opaque through two principal failure modes. First, temporal attenuation: systems that rely on time-to-live (\TTL) constraints or idle timeout policies can destroy agents before their work is confirmed complete, or conversely allow them to persist indefinitely if clocks are misconfigured or communication is delayed. Second, structural mutation: depth-limiting policies constrain how many generational hops may occur, yet they do not prevent descendants from surviving the policy enforcement point and continuing to operate outside the originating system's knowledge horizon. In both cases, the underlying weakness is \textbf{mutability}: the mechanisms that bound agent lifetimes and propagation are external to the agent's identity and can be overwritten, bypassed, or simply forgotten by a crashed parent. Neither TTL counters nor depth counters are part of the agent's immutable operational record. They exist in configuration, not in the execution graph.

These limitations are not merely engineering oversights. They reflect a deeper conceptual gap: existing multi-agent formalisms treat spawning as an ephemeral event rather than as a permanent commitment within a durable lineage. When a parent spawns a child, the child inherits nothing that the system can later use to reconstruct the causal chain. There is no immutable pointer back to the parent, no single canonical identifier that propagates with every descendant, and no structural invariant that prevents a descendant from severing its association with its ancestors. Consequently, the spawn graph is always a mutable approximation—one that reflects the system's current beliefs rather than the ground-truth execution history.

The core insight of this paper is that autonomous drift can be made structurally impossible by a deceptively simple constraint: \textbf{every spawned agent carries an immutable parent pointer that is part of its identity and cannot be removed, spoofed, or lost.} We call this the \textbf{parent pointer chain} (\PPC). When an agent spawns a child, the system creates a new agent whose identity includes a cryptographically signed reference to the parent's own immutable identity record. That record, in turn, references its own parent, and so on, back to the root agent. The chain is append-only: no agent in the chain can be reordered, removed, or assigned a false predecessor. Because every descendant carries its entire ancestry within its identity, the originating system can always reconstruct the complete spawn lineage of any agent, at any time, without relying on external state or mutable configuration.

This invariant transforms drift from a runtime hazard into a structural impossibility. An agent whose parent pointer cannot be excised remains within the system's accountability graph. A descendant cannot survive independently of its chain because there is no mechanism that produces a parentless agent—the system architecture prohibits it by construction. Drift requires the ability to sever the ancestry link; the immutable parent pointer chain removes that ability entirely.

We now provide a formal definition.

\begin{definition}[Autonomous Drift]
Let $G = (V, E)$ be a directed acyclic spawn graph where each vertex $v \in V$ represents an agent and each directed edge $(p, c) \in E$ represents a spawn event from parent $p$ to child $c$. Let $\rho(v)$ denote the root ancestor of $v$ derived by transitively following parent pointers until no parent exists. An agent $v$ is said to be \textbf{autonomously drifted} if and only if the originating system at $t_0$ cannot construct a verifiable path from $\rho(v)$ to $v$ at some later time $t_1 > t_0$, where verifiability requires that every edge along the path is cryptographically signed and immutable.
\end{definition}

This definition distinguishes drift from ordinary agent termination. An agent that completes its task and is gracefully terminated by its parent is not drifted—it exited the active graph through a sanctioned pathway. Drift specifically requires that the ancestry link becomes severed, unavailable, or untrustworthy while the descendant continues to operate. The \PPC invariant guarantees that no such severance is architecturally possible, making drift definitionally unachievable within a compliant system.

The remainder of this paper is organized as follows. Section~\ref{sec:model} presents the formal agent model, including the immutable identity schema and append-only parent pointer construction. Section~\ref{sec:proof} provides a structural proof that the parent pointer chain precludes drift under all spawn scenarios. Section~\ref{sec:implementation} describes the BX3 runtime architecture that enforces the \PPC invariant at the system level, including identity issuance, chain traversal operations, and the role of the Root Agent as the canonical origin of all lineages. Section~\ref{sec:evaluation} presents empirical evaluation of spawn overhead, chain traversal performance, and boundary-condition analysis. Section~\ref{sec:related} surveys related work in multi-agent accountability, spawn control, and lineage-aware computation. Section~\ref{sec:conclusion} offers closing remarks and identifies open extensions for policy-governed drift tolerance and partially immutable chains.